\chapter*{Symbols and Abbreviations}
\addcontentsline{toc}{chapter}{Symbols and Abbreviations}
\markboth{Symbols and Abbreviations}{Symbols and Abbreviations}

The symbols below are defined according to their use in the equations and examples throughout the book. A symbol can carry different meanings in different contexts; where that happens, the accompanying description or equation determines the intended meaning.

\section*{Information Theory and Probability}

\begin{description}
  \item[\(X\)] A random variable modeling the output or message of a source.
  \item[\(x_i\)] One possible outcome or message.
  \item[\(p_i\)] The probability of outcome \(x_i\).
  \item[\(I(x_i)\)] The self-information of outcome \(x_i\), equal to \(-\log_2 p_i\).
  \item[\(H(X)\)] The Shannon entropy of random variable \(X\).
  \item[\(H_2(p)\)] Binary entropy for two outcomes with probabilities \(p\) and \(1-p\).
  \item[\(M\)] The number of possible states or messages.
  \item[\(s\)] The number of symbols available at each position in a string.
  \item[\(n\)] The length of a string or number of its positions.
  \item[\(C\)] Channel capacity; in an electrical context, roman \(\mathrm{C}\) is the unit coulomb.
  \item[\(B\)] The bandwidth of a channel, usually measured in hertz.
  \item[\(S/N\)] The ratio of signal power to noise power.
  \item[\(K_U(x)\)] The Kolmogorov complexity of string \(x\) relative to reference machine \(U\).
  \item[\(U\)] The universal description machine used in the definition of Kolmogorov complexity.
  \item[\(|p|\)] The length of program description \(p\).
\end{description}

\section*{Time, Sampling, and Electricity}

\begin{description}
  \item[\(T_s\)] The time interval between successive samples.
  \item[\(f_s\)] The sampling rate, in samples per second.
  \item[\(f_{\max}\)] The highest frequency present in a band-limited signal.
  \item[\(q\)] Electric charge.
  \item[\(\Delta V\)] The electric potential difference between two points.
  \item[\(\Delta U\)] The change in electric potential energy.
  \item[\(k_{\mathrm B}\)] Boltzmann's constant.
  \item[\(T\)] Absolute temperature, measured in kelvin.
  \item[\(Q_{\min}\)] The lower bound on heat released when one bit is logically erased, within the scope of Landauer's principle.
  \item[\(\ln\)] The natural logarithm, with base \(e\).
  \item[\(\log_2\)] The logarithm with base two.
\end{description}

\section*{Abbreviations}

\begin{description}
  \item[DIKW] \textit{Data, Information, Knowledge, Wisdom}.
  \item[DNA] Deoxyribonucleic acid; the molecule involved in storing and transmitting hereditary information.
  \item[BSC] \textit{Binary Symmetric Channel}.
  \item[SI] The International System of Units.
\end{description}
